\chapter{Complexity — Embed}
%%% generated by the blueprint pipeline from TCSlib.Complexity.TuringMachine.Build.Embed
\section{Overview}

Declarations of the Lean module \texttt{TCSlib.Complexity.TuringMachine.Build.Embed}.

\section{Declarations}

\begin{definition}[Recovering an embedded tape index]
\lean{Turing.embedSlot}
\label{Turing.embedSlot}
\leanok
Given an injective map from the indices of $m$ tapes into the indices of $k$ tapes and
an index $j<k$, return the unique source index mapped to $j$, if one exists, and
otherwise return no index.
\end{definition}

\begin{lemma}[Embedded Tape Index Is Selected]
\lean{Turing.embedSlot_selected}
\label{Turing.embedSlot_selected}
\leanok
\uses{Turing.embedSlot}
\difficulty{4}
Let $m,k$ be natural numbers, let $\iota:\mathrm{Fin}(m)\hookrightarrow\mathrm{Fin}(k)$
be an embedding, and let $i\in\mathrm{Fin}(m)$. The slot selected at tape index
$\iota(i)$ is exactly $i$.
\end{lemma}

\begin{lemma}[Unembedded Tape Index Has No Selected Slot]
\lean{Turing.embedSlot_unselected}
\label{Turing.embedSlot_unselected}
\leanok
\uses{Turing.embedSlot}
\difficulty{3}
Let $m,k$ be natural numbers, let $\iota:\mathrm{Fin}(m)\hookrightarrow\mathrm{Fin}(k)$
be an embedding, and let $j\in\mathrm{Fin}(k)$ lie outside its image. Then no slot is
selected at $j$.
\end{lemma}

\begin{definition}[Embedding a machine action]
\lean{Turing.embedActionCore}
\label{Turing.embedActionCore}
\leanok
\uses{Turing.Action, Turing.embedSlot}
Given an injective map from $m$ work-tape indices into $k$ indices, an optional sink
tape, and an action on $m$ work tapes, define an action on $k$ work tapes. The
input-head movement and successor state are unchanged; mapped tapes retain their
original instructions, while other tapes are left blank and stationary, except that when
a sink tape is specified it writes the original action's output symbol there and moves
right if that output exists. With a sink tape, the expanded action emits no output;
without one, it retains the original output.
\end{definition}

\begin{definition}[Embedding a Silent Configuration]
\lean{Turing.embedSilentCfg}
\label{Turing.embedSilentCfg}
\leanok
\uses{Turing.Cfg, Turing.FinTM.bufferTape, Turing.embedSlot}
Let $m,k$ be nonnegative integers, let $\iota$ be an injective map from the $m$
work-tape indices into the $k$ work-tape indices, and choose a distinguished tape index
$c$ among the $k$ tapes. Given tape contents $T_j:\mathbb
Z\to\mathrm{Option}(\mathrm{Symbol})$, head positions $h_j\in\mathbb Z$, lists $p,o_0$
of symbols, and an $m$-tape configuration $C$ on input $x$, form a $k$-tape
configuration on the same input. It keeps $C$’s state and input-head position, places
each of its work tapes and heads at the index specified by $\iota$, and at every unused
index other than $c$ uses $T_j$ and $h_j$. At index $c$, its tape contains the
concatenation of $p$ and $C$’s output, its head is at the length of that concatenation,
and its output is $o_0$.
\end{definition}

\begin{definition}[Embedding an Emitting Configuration]
\lean{Turing.embedEmitCfg}
\label{Turing.embedEmitCfg}
\leanok
\uses{Turing.Cfg, Turing.embedSlot}
Let $m,k$ be nonnegative integers, let $\iota$ be an injective map from the $m$
work-tape indices into the $k$ work-tape indices, and choose tape contents $T_j:\mathbb
Z\to\mathrm{Option}(\mathrm{Symbol})$ and head positions $h_j\in\mathbb Z$ for the $k$
tapes. Given a symbol list $p$ and an $m$-tape configuration $C$ on input $x$, form a
$k$-tape configuration on the same input. It keeps $C$’s state and input-head position,
places each of its work tapes and heads at the index specified by $\iota$, uses $T_j$
and $h_j$ at every unused index, and has output equal to the concatenation of $p$ and
$C$’s output.
\end{definition}

\begin{definition}[Embedding a Silent Turing Machine]
\lean{Turing.embedSilentTM}
\label{Turing.embedSilentTM}
\leanok
\uses{Turing.MultiTapeTM, Turing.embedActionCore}
Let $m,k$ be nonnegative integers, let $\iota$ inject the $m$ work-tape indices into the
$k$ work-tape indices, and choose one of the $k$ indices as a distinguished tape. From
an $m$-tape Turing machine, form a $k$-tape machine with the same initial state. Its
transition reads the symbols on the tapes selected by $\iota$, applies the original
transition, and embeds the resulting action into the larger tape set using the
distinguished tape for the silent-action convention.
\end{definition}

\begin{definition}[Embedding an Emitting Turing Machine]
\lean{Turing.embedEmitTM}
\label{Turing.embedEmitTM}
\leanok
\uses{Turing.MultiTapeTM, Turing.embedActionCore}
Let $m,k$ be nonnegative integers and let $\iota$ inject the $m$ work-tape indices into
the $k$ work-tape indices. From an $m$-tape Turing machine, form a $k$-tape machine with
the same initial state. Its transition reads the symbols on the tapes selected by
$\iota$, applies the original transition, and embeds the resulting action into the
larger tape set using the emitting-action convention.
\end{definition}

\begin{theorem}[Transition behavior under tape embedding]
\lean{Turing.embedEmitTM_tr}
\label{Turing.embedEmitTM_tr}
\leanok
\uses{Turing.MultiTapeTM, Turing.embedActionCore, Turing.embedEmitTM, Turing.embedSlot_selected, Turing.embedSlot_unselected}
\difficulty{3}
Let $m,k$ be natural numbers, let $S$ be a set of states, let $\text{Symbol}$ be an
alphabet, and let $M$ be a Turing machine with $m$ work tapes. Given an injection from
the $m$ work-tape indices into the $k$ work-tape indices, a state $q\in S$, an optional
input symbol, and a choice of optional symbol for each of the $k$ work tapes, the
transition of the machine obtained by embedding $M$ into $k$ work tapes has the same
input-head movement, output, and successor state as the transition of $M$ on the symbols
at the embedded indices. On each embedded index it has the corresponding work-tape
action of $M$, and on every index outside the image it prescribes no write and zero
movement.
\end{theorem}

\begin{lemma}[Action Commutes with Silent Configuration Embedding]
\lean{Turing.embedSilent_apply}
\label{Turing.embedSilent_apply}
\leanok
\uses{Turing.Action, Turing.Action.apply, Turing.Cfg, Turing.FinTM.bufferTape_append, Turing.embedActionCore, Turing.embedSilentCfg, Turing.embedSlot}
\difficulty{5}
Let $m,k$ be natural numbers, let $\iota:\mathrm{Fin}(m)\hookrightarrow\mathrm{Fin}(k)$,
and choose a tape index $\mathrm{cap}$. For any $m$-tape action $a$, any initial
configuration $c$, any $k$-tape contents and head positions, and any words
$\mathrm{pre}$ and $\mathrm{out}_0$, applying the embedded action to the embedded
configuration gives exactly the embedded configuration obtained by first applying $a$ to
$c$.
\end{lemma}

\begin{lemma}[One Step Commutes with Silent Machine Embedding]
\lean{Turing.embedSilent_step}
\label{Turing.embedSilent_step}
\leanok
\uses{Turing.Action.apply, Turing.Cfg, Turing.Cfg.inputSymbol, Turing.Cfg.workTapeSymbols, Turing.MultiTapeTM, Turing.MultiTapeTM.step, Turing.embedActionCore, Turing.embedSilentCfg, Turing.embedSilentTM, Turing.embedSilent_apply, Turing.embedSlot_selected}
\difficulty{4}
Let $m,k$ be natural numbers, let $\iota:\mathrm{Fin}(m)\hookrightarrow\mathrm{Fin}(k)$,
and choose a tape index $\mathrm{cap}$. For any $m$-tape machine $M$, any initial
configuration $c$, any $k$-tape contents and head positions, and any words
$\mathrm{pre}$ and $\mathrm{out}_0$, one step of the silently embedded machine from the
embedded configuration equals the embedding of the configuration obtained by one step of
$M$ from $c$.
\end{lemma}

\begin{theorem}[Run commutes with silent embedding]
\lean{Turing.embedSilentTM_runFrom}
\label{Turing.embedSilentTM_runFrom}
\leanok
\uses{Turing.Cfg, Turing.MultiTapeTM, Turing.MultiTapeTM.runFrom, Turing.MultiTapeTM.runFrom_comm_of_step, Turing.embedSilentCfg, Turing.embedSilentTM, Turing.embedSilent_step}
\difficulty{1}
Let $\iota$ embed the $m$ tape indices into the $k$ tape indices, and let the
distinguished index $\mathrm{cap}$ lie outside its image. For any $m$-tape machine $M$,
initial configuration $c$, background tape contents and head positions, prefix and
initial output, and natural number $t$, running the silent embedding for $t$ steps from
the embedded configuration gives exactly the embedding of the configuration reached by
running $M$ from $c$ for $t$ steps.
\end{theorem}

\begin{theorem}[Frame properties of the silent embedding]
\lean{Turing.embedSilentTM_frame}
\label{Turing.embedSilentTM_frame}
\leanok
\uses{Turing.Cfg, Turing.MultiTapeTM, Turing.MultiTapeTM.runFrom, Turing.embedSilentCfg, Turing.embedSilentTM, Turing.embedSilentTM_runFrom, Turing.embedSlot_unselected}
\difficulty{3}
Let $\iota$ embed the $m$ tape indices into the $k$ tape indices, with distinguished
index $\mathrm{cap}$ outside its image, and let $M$ be an $m$-tape machine. For any
initial configuration, background tape contents and head positions, prefix, initial
output, and time $t$, every tape index outside the image of $\iota$ and different from
$\mathrm{cap}$ retains its initial tape contents and head position after $t$ steps of
the silent embedding. The input-head position agrees with that of $M$ after $t$ steps,
and the output remains the initial output.
\end{theorem}

\begin{theorem}[Visited positions on embedded tapes]
\lean{Turing.embedSilentTM_visitedByTapeHead}
\label{Turing.embedSilentTM_visitedByTapeHead}
\leanok
\uses{Turing.Cfg, Turing.MultiTapeTM, Turing.MultiTapeTM.spaceUsedByTape, Turing.MultiTapeTM.visitedByTapeHead, Turing.embedSilentCfg, Turing.embedSilentTM, Turing.embedSilentTM_runFrom, Turing.embedSlot_selected}
\difficulty{4}
Let $\iota$ embed the $m$ tape indices into the $k$ tape indices, with distinguished
index $\mathrm{cap}$ outside its image, and let $M$ be an $m$-tape machine. For any
initial configuration, background tape contents and head positions, prefix, initial
output, natural number $t$, and tape index $i<m$, the set of head positions visited
through time $t$ on tape $\iota(i)$ by the silent embedding equals the set visited on
tape $i$ by $M$. Consequently, the numbers of distinct positions in these two sets are
equal.
\end{theorem}

\begin{theorem}[Unmapped Tape Head Remains Fixed]
\lean{Turing.embedSilentTM_visitedByTapeHead_frame}
\label{Turing.embedSilentTM_visitedByTapeHead_frame}
\leanok
\uses{Turing.Cfg, Turing.MultiTapeTM, Turing.MultiTapeTM.spaceUsedByTape, Turing.MultiTapeTM.visitedByTapeHead, Turing.embedSilentCfg, Turing.embedSilentTM, Turing.embedSilentTM_runFrom, Turing.embedSlot_unselected}
\difficulty{4}
Let $m,k$ be natural numbers, let $\iota:\mathrm{Fin}(m)\hookrightarrow\mathrm{Fin}(k)$
be an embedding, and choose a tape index $\mathrm{cap}$ outside its image. Let $M$ be an
$m$-work-tape machine, and let $\mathrm{tapes}$ and $\mathrm{heads}$ specify the
contents and head positions of the $k$ tapes; let $\mathrm{pre}$ and $\mathrm{out}_0$ be
words and $c$ an initial configuration of $M$. For any time $t$ and tape index $j$
outside the image of $\iota$ with $j\ne\mathrm{cap}$, the head of tape $j$ in the
embedded computation visits exactly the singleton $\{\mathrm{heads}(j)\}$ through time
$t$, and that tape uses exactly one cell.
\end{theorem}

\begin{theorem}[Space bound for the output tape]
\lean{Turing.embedSilentTM_spaceUsedByTape_cap}
\label{Turing.embedSilentTM_spaceUsedByTape_cap}
\leanok
\uses{Turing.Cfg, Turing.MultiTapeTM, Turing.MultiTapeTM.output_prefix, Turing.MultiTapeTM.runFrom, Turing.MultiTapeTM.spaceUsedByTape, Turing.MultiTapeTM.visitedByTapeHead, Turing.embedSilentCfg, Turing.embedSilentTM, Turing.embedSilentTM_runFrom, Turing.embedSlot_unselected}
\difficulty{5}
Let $\iota$ embed the $m$ tape indices into the $k$ tape indices, with distinguished
index $\mathrm{cap}$ outside its image, and let $M$ be an $m$-tape machine. For any
initial configuration $c$, background tape contents and head positions, prefix and
initial output, and time $t$, the number of distinct positions visited by the head of
tape $\mathrm{cap}$ in the silent embedding through time $t$ is at most the length of
$M$'s output after $t$ steps minus the length of its initial output, plus one.
\end{theorem}

\begin{lemma}[Action on an embedded configuration]
\lean{Turing.embedEmit_apply}
\label{Turing.embedEmit_apply}
\leanok
\uses{Turing.Action, Turing.Action.apply, Turing.Cfg, Turing.embedActionCore, Turing.embedEmitCfg, Turing.embedSlot}
\difficulty{4}
Let $m,k$ be natural numbers, let $\iota$ embed the $m$ work-tape indices into the $k$
work-tape indices, and choose arbitrary initial contents and head positions for the $k$
tapes, an output prefix, an initial configuration on $m$ tapes, and a one-step action on
those tapes. Applying the action after embedding the configuration—with output emitted
on the output tape—gives the same configuration as first applying the action to the
original configuration and then embedding the result with the same initial contents,
head positions, and output prefix.
\end{lemma}

\begin{lemma}[Step of an embedded machine]
\lean{Turing.embedEmit_step}
\label{Turing.embedEmit_step}
\leanok
\uses{Turing.Action.apply, Turing.Cfg, Turing.Cfg.inputSymbol, Turing.Cfg.workTapeSymbols, Turing.MultiTapeTM, Turing.MultiTapeTM.step, Turing.embedActionCore, Turing.embedEmitCfg, Turing.embedEmitTM, Turing.embedEmit_apply, Turing.embedSlot_selected}
\difficulty{5}
Let $m,k$ be natural numbers, let $\iota$ be an embedding of the $m$ work-tape indices
into the $k$ work-tape indices, and let $M$ be a Turing machine with $m$ work tapes. For
any $k$-tape contents and head positions, any output prefix, and any configuration $c$
of $M$, taking one step of the machine obtained by embedding $M$ into $k$ work tapes
from the configuration formed by embedding $c$ with those contents, head positions, and
output prefix gives the configuration formed by embedding the result of one step of $M$
from $c$ with the same contents, head positions, and output prefix.
\end{lemma}

\begin{theorem}[Run commutes with tape embedding]
\lean{Turing.embedEmitTM_runFrom}
\label{Turing.embedEmitTM_runFrom}
\leanok
\uses{Turing.Cfg, Turing.MultiTapeTM, Turing.MultiTapeTM.runFrom, Turing.MultiTapeTM.runFrom_comm_of_step, Turing.embedEmitCfg, Turing.embedEmitTM, Turing.embedEmit_step}
\difficulty{1}
Let $M$ be a machine with $m$ work tapes, embed those tape indices into $k$ work tapes,
and form a $k$-tape configuration by placing $c$'s tapes and head positions at the
embedded indices, using arbitrary initial contents and head positions elsewhere, and
prefixing its output by a fixed word. For every natural number $t$, running the embedded
machine for $t$ steps gives exactly the configuration formed in the same way from the
configuration reached by running $M$ from $c$ for $t$ steps.
\end{theorem}

\begin{theorem}[Frame preservation under tape embedding]
\lean{Turing.embedEmitTM_frame}
\label{Turing.embedEmitTM_frame}
\leanok
\uses{Turing.Cfg, Turing.MultiTapeTM, Turing.MultiTapeTM.runFrom, Turing.embedEmitCfg, Turing.embedEmitTM, Turing.embedEmitTM_runFrom, Turing.embedSlot_unselected}
\difficulty{3}
Let $x$ be an input word, let $M$ be a deterministic machine with $m$ work tapes, and
let $\iota : \{0,\ldots,m-1\} \hookrightarrow \{0,\ldots,k-1\}$ be an embedding of tape
indices. Choose initial contents $T_j : \mathbb{Z} \to \mathrm{Option}(\mathrm{Symbol})$
and head positions $h_j \in \mathbb{Z}$ for each of the $k$ tapes, an output prefix $p$,
an $m$-tape configuration $c$ on input $x$, and $t \in \mathbb{N}$. After $t$ steps of
the $k$-tape machine obtained from $M$, starting with $M$'s tapes embedded by $\iota$,
the other tapes still have contents $T_j$ and head positions $h_j$; the input-head
position agrees with the position after $t$ steps of $M$ from $c$, and the output is $p$
followed by the output after those $t$ steps of $M$.
\end{theorem}

\begin{theorem}[Visited positions and space on embedded tapes]
\lean{Turing.embedEmitTM_visitedByTapeHead}
\label{Turing.embedEmitTM_visitedByTapeHead}
\leanok
\uses{Turing.Cfg, Turing.MultiTapeTM, Turing.MultiTapeTM.spaceUsedByTape, Turing.MultiTapeTM.visitedByTapeHead, Turing.embedEmitCfg, Turing.embedEmitTM, Turing.embedEmitTM_runFrom, Turing.embedSlot_selected}
\difficulty{4}
Let $M$ be a machine with $m$ work tapes, embed those tape indices into $k$ work-tape
indices, and form a $k$-tape configuration from an $m$-tape configuration $c$ using
arbitrary initial contents and head positions on the other tapes and an output prefix.
For every time $t$ and every original tape index $i$, the set of positions visited by
head $i$ in $M$ through time $t$ equals the set visited by head $\sigma(i)$ in the
embedded simulation through time $t$; consequently, the numbers of distinct positions in
those sets are equal.
\end{theorem}

\begin{theorem}[Unvisited embedded work tape]
\lean{Turing.embedEmitTM_visitedByTapeHead_frame}
\label{Turing.embedEmitTM_visitedByTapeHead_frame}
\leanok
\uses{Turing.Cfg, Turing.MultiTapeTM, Turing.MultiTapeTM.spaceUsedByTape, Turing.MultiTapeTM.visitedByTapeHead, Turing.embedEmitCfg, Turing.embedEmitTM, Turing.embedEmitTM_runFrom, Turing.embedSlot_unselected}
\difficulty{4}
Let $m,k$ be natural numbers, let $\iota$ embed the $m$ work-tape indices into the $k$
work-tape indices, and let $M$ be an $m$-work-tape Turing machine. For any initial tape
contents and head positions, output prefix, initial configuration, time $t$, and index
$j$ outside the image of $\iota$, the set of positions visited by tape $j$ in the
embedded computation through time $t$ is the singleton consisting of its initial head
position; consequently, that tape uses exactly one cell.
\end{theorem}

\begin{definition}[Embedding a Silent Machine with Return State]
\lean{Turing.embedSilentRetTM}
\label{Turing.embedSilentRetTM}
\leanok
\uses{Turing.MultiTapeTM, Turing.embedActionCore}
Let $m,k$ be nonnegative integers, let $\iota$ inject the $m$ work-tape indices into the
$k$ work-tape indices, and choose one of the $k$ indices as a distinguished tape. From
an $m$-tape Turing machine with state set $S$, form a $k$-tape machine with states
$S\sqcup\{\ast\}$, where $\ast$ is a new state. Its initial state is the original
initial state. In an original state, it reads the symbols on tapes selected by $\iota$,
applies the original transition, embeds that action using the distinguished tape, and
sends a halting transition to $\ast$ instead of halting. In state $\ast$, it stays on
the input, makes no work-tape writes or head movements, emits no symbol, and transitions
to $\ast$.
\end{definition}

\begin{definition}[Embedding an Emitting Machine with Return State]
\lean{Turing.embedEmitRetTM}
\label{Turing.embedEmitRetTM}
\leanok
\uses{Turing.MultiTapeTM, Turing.embedActionCore}
Let $m,k$ be nonnegative integers and let $\iota$ inject the $m$ work-tape indices into
the $k$ work-tape indices. From an $m$-tape Turing machine with state set $S$, form a
$k$-tape machine with states $S\sqcup\{\ast\}$, where $\ast$ is a new state. Its initial
state is the original initial state. In an original state, it reads the symbols on tapes
selected by $\iota$, applies the original transition, embeds that action using the
emitting-action convention, and sends a halting transition to $\ast$ instead of halting.
In state $\ast$, it stays on the input, makes no work-tape writes or head movements,
emits no symbol, and transitions to $\ast$.
\end{definition}

\begin{definition}[Action with a Return State]
\lean{Turing.embedReturnAction}
\label{Turing.embedReturnAction}
\leanok
\uses{Turing.Action}
Let $a$ be a one-step action for a machine with state set $S$. Regard it as an action
with state set $S\sqcup\{\ast\}$ by retaining its input-head movement, work-tape
actions, and output, and replacing its successor state with the original successor when
present or with the new state $\ast$ when the original action halts.
\end{definition}

\begin{definition}[Configuration with a Return State]
\lean{Turing.embedReturnCfg}
\label{Turing.embedReturnCfg}
\leanok
\uses{Turing.Cfg}
Let $C$ be a configuration on input $x$ with state set $S$. Regard it as a configuration
with state set $S\sqcup\{\ast\}$ by retaining its input-head position, work tapes,
work-tape head positions, and output, and replacing its state with the original state
when present or with the new state $\ast$ when it has halted.
\end{definition}

\begin{lemma}[Live state under return embedding]
\lean{Turing.embedReturnCfg_live}
\label{Turing.embedReturnCfg_live}
\leanok
\uses{Turing.Cfg, Turing.Cfg.mapState, Turing.embedReturnCfg}
\difficulty{4}
Let $k$ be a natural number and let $c$ be a configuration of a $k$-work-tape machine.
If $c$ has a state, then embedding its state into the left summand of the disjoint union
of its state set with a singleton gives the same configuration as renaming its state by
the left-summand injection.
\end{lemma}

\begin{lemma}[Step under return embedding]
\lean{Turing.embedReturn_step}
\label{Turing.embedReturn_step}
\leanok
\uses{Turing.Action.apply, Turing.Cfg, Turing.Cfg.inputSymbol, Turing.Cfg.workTapeSymbols, Turing.MultiTapeTM, Turing.MultiTapeTM.step, Turing.embedReturnAction, Turing.embedReturnCfg}
\difficulty{4}
Let $k$ be a natural number, let $N$ be a $k$-work-tape machine with state set $S$, and
let $R$ be such a machine with state set $S \sqcup \{\ast\}$. Suppose that, at every
state in the left summand, $R$'s transition is the return embedding of $N$'s transition,
and that at the added state its transition makes no tape changes or output, stays in
that state, and does not move either the input head or any work-tape head. Then, for
every configuration $c$ of $N$, taking one step of $R$ from the return embedding of $c$
gives the return embedding of the configuration obtained by taking one step of $N$.
\end{lemma}

\begin{lemma}[One step of the silent return machine]
\lean{Turing.embedSilentRet_step}
\label{Turing.embedSilentRet_step}
\leanok
\uses{Turing.Cfg, Turing.Cfg.mapState, Turing.MultiTapeTM, Turing.MultiTapeTM.step, Turing.embedReturnCfg, Turing.embedReturnCfg_live, Turing.embedReturn_step, Turing.embedSilentCfg, Turing.embedSilentRetTM, Turing.embedSilentTM, Turing.embedSilent_step}
\difficulty{3}
Let $\iota$ embed the $m$ tape indices into the $k$ tape indices, let $\mathrm{cap}$ be
a distinguished $k$-tape index, and let $M$ be an $m$-tape machine. For any initial tape
contents and head positions, prefix and initial output, and configuration $c$ whose
state is present, one step of the return machine from the embedded configuration with
state mapped into the left summand equals the return-state embedding of the
configuration reached by one step of $M$ from $c$.
\end{lemma}

\begin{lemma}[Simulation Through the Halting Step]
\lean{Turing.embedThroughHalt}
\label{Turing.embedThroughHalt}
\leanok
\uses{Turing.Cfg, Turing.Cfg.mapState, Turing.MultiTapeTM, Turing.MultiTapeTM.runFrom, Turing.MultiTapeTM.runFrom_succ_eq_step', Turing.MultiTapeTM.step, Turing.embedReturnCfg, Turing.embedReturnCfg_live}
\difficulty{5}
Let $M$ be an $m$-tape machine and $R$ a $k$-tape machine whose states are in
$S\oplus\mathrm{Unit}$. Suppose $E$ maps configurations of $M$ to configurations with
state set $S$ and preserves their optional state components. Assume that whenever a
configuration $d$ of $M$ is not halted, one step of $R$ from the state-renamed image of
$d$ is the return encoding of $E$ applied to the result of one step of $M$. Let $c$ be a
nonhalted initial configuration and let $T$ be a natural number such that every
configuration of $M$ at times $t<T$ is nonhalted, while its configuration at time $T$ is
halted. Then for every $t<T$, the run of $R$ from the state-renamed image of $E(c)$
equals the state-renamed image of $E$ of the configuration of $M$ at time $t$; at time
$T$ it equals the configuration obtained from $E$ of $M$'s time-$T$ configuration by
setting its state to the right-hand summand; and at every time $t<T$, its state is not
that right-hand summand.
\end{lemma}

\begin{theorem}[Run of the silent return machine]
\lean{Turing.embedSilentRetTM_run}
\label{Turing.embedSilentRetTM_run}
\leanok
\uses{Turing.Cfg, Turing.Cfg.mapState, Turing.MultiTapeTM, Turing.MultiTapeTM.runFrom, Turing.embedSilentCfg, Turing.embedSilentRetTM, Turing.embedSilentRet_step, Turing.embedThroughHalt}
\difficulty{1}
Let $\iota$ embed the $m$ tape indices into the $k$ tape indices, and let the
distinguished tape index $\mathrm{cap}$ lie outside the image of $\iota$. For an
$m$-tape machine $M$ and configuration $c$ whose state is present, suppose that $M$ is
still in a state at every time $t<T$ and has halted at time $T$. Starting from the
embedded configuration with its state mapped into the left summand, the return machine's
run agrees, at every time $t<T$, with the embedding of $M$'s configuration at time $t$
(again with state in the left summand); at time $T$ it is the embedded halted
configuration with state set to the right-summand return state. At every time before
$T$, its state is not that return state.
\end{theorem}

\begin{lemma}[One-step return-state simulation]
\lean{Turing.embedEmitRet_step}
\label{Turing.embedEmitRet_step}
\leanok
\uses{Turing.Cfg, Turing.Cfg.mapState, Turing.MultiTapeTM, Turing.MultiTapeTM.step, Turing.embedEmitCfg, Turing.embedEmitRetTM, Turing.embedEmitTM, Turing.embedEmit_step, Turing.embedReturnCfg, Turing.embedReturnCfg_live, Turing.embedReturn_step}
\difficulty{3}
Let $M$ be a machine with $m$ work tapes, embed those tape indices into $k$ work tapes,
and let $c$ be an $m$-tape configuration whose state is not halted. For any initial
contents and head positions on the other tapes and any output prefix, one step of the
simulation with a distinguished return state, from the embedded configuration with
states renamed into the left summand, equals the embedded configuration obtained after
one step of $M$, with its state changed to the left summand when it remains active and
to the distinguished right-summand state when it halts.
\end{lemma}

\begin{theorem}[Run of the return-state simulation]
\lean{Turing.embedEmitRetTM_run}
\label{Turing.embedEmitRetTM_run}
\leanok
\uses{Turing.Cfg, Turing.Cfg.mapState, Turing.MultiTapeTM, Turing.MultiTapeTM.runFrom, Turing.embedEmitCfg, Turing.embedEmitRetTM, Turing.embedEmitRet_step, Turing.embedThroughHalt}
\difficulty{1}
Let $M$ be a machine with $m$ work tapes, and embed those tape indices into $k$ work
tapes. Start its $k$-tape simulation from the configuration obtained by placing its
tapes at the embedded indices, supplying arbitrary contents and head positions on the
other tapes, prefixing its output by a fixed word, and renaming each original state into
a left-summand state. Suppose the initial configuration is not halted, every
configuration reached by $M$ at a time $t<T$ is not halted, and the configuration at
time $T$ is halted. Then, at every time $t<T$, the simulation agrees with the embedded
configuration of $M$ with its state renamed into the left summand; at time $T$ it agrees
in all configuration data except that its state is the distinguished right-summand
state; and at all times $t<T$ its state is not that distinguished state.
\end{theorem}

\begin{lemma}[Visited cells under return embedding]
\lean{Turing.embedReturn_visited}
\label{Turing.embedReturn_visited}
\leanok
\uses{Turing.Cfg, Turing.Cfg.mapState, Turing.MultiTapeTM, Turing.MultiTapeTM.runFrom_comm_of_step, Turing.MultiTapeTM.runFrom_of_halt, Turing.MultiTapeTM.visitedByTapeHead, Turing.embedReturnAction, Turing.embedReturnCfg, Turing.embedReturnCfg_live, Turing.embedReturn_step}
\difficulty{4}
Let $k$ be a natural number, let $N$ be a $k$-work-tape machine with state set $S$, and
let $R$ be such a machine with state set $S \sqcup \{\ast\}$. Suppose that at every
left-summand state, $R$'s transition is the return embedding of $N$'s transition, and
that at the added state its transition makes no tape changes or output, stays in that
state, and does not move either the input head or any work-tape head. For every
configuration $c$ of $N$, time $t$, and work-tape index $j$, the set of positions
visited by tape $j$ through time $t$ by $R$ from the configuration obtained by renaming
$c$'s state into the left summand equals the set visited by tape $j$ through time $t$ by
$N$ from $c$.
\end{lemma}

\begin{theorem}[Visited head positions on embedded tapes]
\lean{Turing.embedSilentRetTM_visitedByTapeHead}
\label{Turing.embedSilentRetTM_visitedByTapeHead}
\leanok
\uses{Turing.Cfg, Turing.Cfg.mapState, Turing.MultiTapeTM, Turing.MultiTapeTM.visitedByTapeHead, Turing.embedReturn_visited, Turing.embedSilentCfg, Turing.embedSilentRetTM, Turing.embedSilentTM}
\difficulty{1}
Let $\iota$ embed the $m$ tape indices into the $k$ tape indices, and let $M$ be an
$m$-tape machine. For any natural number $t$, any $k$-tape index $j$, and any initial
configuration and embedding data, the set of head positions visited through time $t$ on
tape $j$ by the return machine, starting from the embedded configuration with its state
mapped into the left summand, equals the set visited on tape $j$ by the silent embedding
of $M$ from the embedded configuration.
\end{theorem}

\begin{theorem}[Visited positions under return-state simulation]
\lean{Turing.embedEmitRetTM_visitedByTapeHead}
\label{Turing.embedEmitRetTM_visitedByTapeHead}
\leanok
\uses{Turing.Cfg, Turing.Cfg.mapState, Turing.MultiTapeTM, Turing.MultiTapeTM.visitedByTapeHead, Turing.embedEmitCfg, Turing.embedEmitRetTM, Turing.embedEmitTM, Turing.embedReturn_visited}
\difficulty{1}
Let $M$ be a machine with $m$ work tapes, and embed those tape indices into $k$ work
tapes. For any initial contents and head positions on the $k$ tapes, output prefix,
$m$-tape configuration $c$, time $t$, and $k$-tape index $j$, the set of positions
visited by tape head $j$ through time $t$ in the simulation with a distinguished return
state, started from the embedded configuration with states renamed into the left
summand, equals the corresponding set for the simulation without the return state
started from the embedded configuration.
\end{theorem}

\begin{theorem}[Selected tape under configuration embedding]
\lean{Turing.embedSilentCfg_selected_tape}
\label{Turing.embedSilentCfg_selected_tape}
\leanok
\uses{Turing.Cfg, Turing.embedSilentCfg, Turing.embedSlot_selected}
\difficulty{1}
Let $m,k$ be tape counts, let $\iota$ embed the $m$ tape indices into the $k$ tape
indices, and let $c$ be an $m$-tape configuration. For every $i<m$, the content of the
$k$-tape configuration obtained by embedding $c$ on tape $\iota(i)$ equals the content
of tape $i$ in $c$.
\end{theorem}

\begin{theorem}[Head position on an embedded tape]
\lean{Turing.embedSilentCfg_selected_pos}
\label{Turing.embedSilentCfg_selected_pos}
\leanok
\uses{Turing.Cfg, Turing.embedSilentCfg, Turing.embedSlot_selected}
\difficulty{1}
Let $m,k$ be natural numbers, let $\iota$ embed the $m$ work-tape indices into the $k$
work-tape indices, and choose a tape index $\mathrm{cap}$, arbitrary initial contents
and head positions, lists $\mathrm{pre}$ and $\mathrm{out}_0$, an initial configuration
$c$, and an index $i$ among the original $m$ tapes. In the embedded silent
configuration, the head position on tape $\iota(i)$ equals its position on tape $i$ in
$c$.
\end{theorem}

\begin{theorem}[Preservation of selected tape contents]
\lean{Turing.embedEmitCfg_selected_tape}
\label{Turing.embedEmitCfg_selected_tape}
\leanok
\uses{Turing.Cfg, Turing.embedEmitCfg, Turing.embedSlot_selected}
\difficulty{1}
Let $m,k$ be nonnegative integers, let $\iota$ embed the indices of $m$ tapes into those
of $k$ tapes, and let $c$ be an $m$-tape configuration. After embedding $c$ into $k$
tapes with any supplied tape contents, head positions, and output prefix, the tape
contents at every integer cell of tape $\iota(i)$ are exactly those of tape $i$ in $c$,
for every $i<m$.
\end{theorem}

\begin{theorem}[Preservation of selected head positions]
\lean{Turing.embedEmitCfg_selected_pos}
\label{Turing.embedEmitCfg_selected_pos}
\leanok
\uses{Turing.Cfg, Turing.embedEmitCfg, Turing.embedSlot_selected}
\difficulty{1}
Let $m,k$ be nonnegative integers, let $\iota$ embed the indices of $m$ tapes into those
of $k$ tapes, and let $c$ be an $m$-tape configuration. After embedding $c$ into $k$
tapes with any supplied tape contents, head positions, and output prefix, the head
position on the tape indexed by $\iota(i)$ equals the head position on tape $i$ in $c$,
for every $i<m$.
\end{theorem}

\begin{theorem}[Contents and head position of the buffer tape]
\lean{Turing.embedSilentCfg_cap}
\label{Turing.embedSilentCfg_cap}
\leanok
\uses{Turing.Cfg, Turing.FinTM.bufferTape, Turing.embedSilentCfg, Turing.embedSlot_unselected}
\difficulty{1}
Let $m,k$ be natural numbers, let $\iota$ embed the $m$ work-tape indices into the $k$
work-tape indices, and choose a tape index $\mathrm{cap}$ outside the image of $\iota$.
For any initial tape contents and head positions, lists $\mathrm{pre}$ and
$\mathrm{out}_0$, and an initial configuration $c$, in the embedded silent configuration
the contents of tape $\mathrm{cap}$ are the buffer storing $\mathrm{pre}$ followed by
$c$'s output, and its head is at the integer position equal to the length of that
concatenated list.
\end{theorem}

\begin{theorem}[Unselected tape preservation]
\lean{Turing.embedEmitCfg_unselected}
\label{Turing.embedEmitCfg_unselected}
\leanok
\uses{Turing.Cfg, Turing.MultiTapeTM, Turing.MultiTapeTM.runFrom_zero, Turing.embedEmitCfg, Turing.embedEmitTM_frame}
\difficulty{3}
Let $m,k$ be natural numbers, let $\sigma$ embed the $m$ work-tape indices into the $k$
work-tape indices, and let $c$ be a configuration on $m$ tapes. For any initial contents
and head positions on $k$ tapes, any output prefix, and any index $j$ outside the image
of $\sigma$, the embedded configuration has tape $j$ equal to its given initial tape and
head position equal to its given initial head position.
\end{theorem}

\begin{definition}[Selection of work tapes]
\lean{Turing.Cfg.select}
\label{Turing.Cfg.select}
\leanok
\uses{Turing.Cfg}
Let $d$ be a configuration of a $k$-work-tape Turing machine, and let
$\sigma:\{0,\ldots,m-1\}\to\{0,\ldots,k-1\}$ be an injection selecting $m$ of its work
tapes. The selected configuration has the same state and input-head position as $d$,
uses the selected tapes and their head positions in the order specified by $\sigma$, and
has an empty output tape.
\end{definition}

\begin{theorem}[Reconstruction of an embedded configuration]
\lean{Turing.embedEmitCfg_eq}
\label{Turing.embedEmitCfg_eq}
\leanok
\uses{Turing.Cfg, Turing.embedEmitCfg, Turing.embedEmitCfg_selected_pos, Turing.embedEmitCfg_selected_tape, Turing.embedEmitCfg_unselected}
\difficulty{4}
Let $m,k$ be nonnegative integers, let $\sigma$ embed the indices of $m$ tapes into
those of $k$ tapes, and let $c$ and $d$ be configurations on the same input, with $c$
using state set $S$ and $d$ using the same state set and symbol set. Suppose their
states and input-head positions agree, and that concatenating a word $p$ with $c$'s
output gives $d$'s output. Suppose also that on every selected tape, $c$'s tape contents
and head position agree with those of $d$; on every unselected tape, the supplied tape
contents and head position agree with those of $d$. Then embedding $c$ using the
supplied tapes and head positions and prefixing its output by $p$ yields $d$.
\end{theorem}

\begin{theorem}[State-renamed configuration reconstruction]
\lean{Turing.embedEmitCfg_mapState_eq}
\label{Turing.embedEmitCfg_mapState_eq}
\leanok
\uses{Turing.Cfg, Turing.Cfg.mapState, Turing.embedEmitCfg, Turing.embedEmitCfg_eq}
\difficulty{1}
Let $m,k$ be nonnegative integers, let $\sigma$ embed the indices of $m$ tapes into
those of $k$ tapes, and let $f:S\to H$ be a map of state sets. Let $c$ be a
configuration with state set $S$ and let $d$ be a configuration with state set $H$, both
on the same input and over the same symbol set. Suppose applying $f$ to the optional
state of $c$ gives the state of $d$, their input-head positions agree, and prefixing
$c$'s output by a word $p$ gives $d$'s output. Suppose their tape contents and head
positions agree on selected tapes, and the supplied tape contents and head positions
agree with $d$ on every unselected tape. Then embedding $c$ with those supplied tapes
and head positions and prefixing its output by $p$, followed by renaming its state using
$f$, gives $d$.
\end{theorem}

\begin{theorem}[Embedding configurations formed from tape words]
\lean{Turing.embedEmitCfg_ofWords_eq}
\label{Turing.embedEmitCfg_ofWords_eq}
\leanok
\uses{Turing.Cfg.ofWords, Turing.FinTM.bufferTape, Turing.embedEmitCfg, Turing.embedEmitCfg_eq}
\difficulty{2}
Let $m,k$ be nonnegative integers, let $\sigma$ embed the indices of $m$ tapes into
those of $k$ tapes, and let $q$ be a state. For each of the $m$ tapes, let $v_i$ be a
Boolean word; for each of the $k$ tapes, let $\mathrm{frame}_j$ and $\mathrm{target}_j$
be Boolean words. Suppose $\mathrm{target}_{\sigma(i)}=v_i$ on every selected tape and
$\mathrm{frame}_j=\mathrm{target}_j$ on every unselected tape. Then embedding the
configuration formed from the words $v_i$, using the framed words on unselected tapes
and an empty output prefix, gives the configuration formed from the target words, on the
same Boolean input and in state $q$.
\end{theorem}

\begin{theorem}[Selecting and re-embedding a configuration]
\lean{Turing.embedEmitCfg_select}
\label{Turing.embedEmitCfg_select}
\leanok
\uses{Turing.Cfg, Turing.Cfg.select, Turing.embedEmitCfg, Turing.embedEmitCfg_eq}
\difficulty{2}
Let $m$ and $k$ be nonnegative integers, let $S$ and $\Gamma$ be types, let $x$ be a
finite input string over $\Gamma$, and let $\sigma$ be an injective map from the $m$
tape indices to the $k$ tape indices. For every $k$-tape configuration $d$ on input $x$,
selecting the tapes indexed by $\sigma$ to form an $m$-tape configuration with the same
state and input-head position and empty output, then embedding it back using $d$'s
work-tape contents and head positions on the other tapes and $d$'s output as the prefix,
recovers $d$ exactly.
\end{theorem}
